% GENERATED FILE. Edit canonical lessons, then run npm run generate:books.
% canonical-source-hash: fnv1a64:1ca6cc8e0ebf5562
% canonical-lessons: FR-C297-carambolage, FR-C297-chauffard, FR-C297-exces-de-vitesse, FR-C297-constat, FR-C297-derapage

\chapter{Accident de la route}
\label{ch:fr-accident-de-la-route}
\hlchaptermodality{297}

\begin{chapteropening}
\textbf{By the end of this chapter:} \emph{I can describe a road accident: a pile-up, a reckless driver, speeding, the accident report and a skid.}

Complete the last lesson of chapter 297: I can describe a road accident: a pile-up, a reckless driver, speeding, the accident report and a skid.
\end{chapteropening}

\section[le carambolage]{le carambolage — the pile-up, the multiple crash}

\label{lesson:FR-C297-carambolage}



\begin{quote}
Before the new one: say the French for the ransom, then the French for to escape.
\end{quote}

\subsection*{You'll want to know: le carambolage}
\textbf{le carambolage} — \textquotedblleft{}the pile-up, the multiple crash\textquotedblright{}.

\begin{quote}
Il y a eu un carambolage sur l'autoroute ce matin.
\end{quote}

\begin{quote}
\emph{There was a pile-up on the motorway this morning.}
\end{quote}

A billiards word first: \emph{caramboler} is to hit two balls in one shot.

\begin{cousinweb}
\emph{carambolage} is built on \emph{caramboler} \textquotedblleft{}to hit two balls at once\textquotedblright{}, from \emph{carambole}, the red ball in billiards, which was borrowed from Spanish \emph{carambola} (EtymDB lists the noun without a parent).
\end{cousinweb}

\begin{grammarlens}[title={What you've built}]
One more word for telling, explaining and describing at B1.
\end{grammarlens}

\subsection*{Your turn}
\begin{itemize}
\raggedright
  \item \emph{Say it:} \emph{le carambolage}
  \item \emph{Say it:} \emph{le carambolage}, once more
  \item \emph{Say it:} a sentence of your own with \emph{le carambolage}
  \item \emph{Recall:} say the French for the ransom, then the French for to escape, then say \emph{le carambolage} again
\end{itemize}

\begin{tcolorbox}[breakable,colback=teal!4,colframe=teal!35!black,title={Before you move on}]
What does le carambolage mean? (\textquotedblleft{}The pile-up, the multiple crash\textquotedblright{}.) Say it once more.
\end{tcolorbox}
\section[le chauffard]{le chauffard — the reckless driver, the road hog}

\label{lesson:FR-C297-chauffard}



\begin{quote}
Before the new one: say the French for to escape, then the French for the pile-up.
\end{quote}

\subsection*{You'll want to know: le chauffard}
\textbf{le chauffard} — \textquotedblleft{}the reckless driver, the road hog\textquotedblright{}.

\begin{quote}
Un chauffard a grillé le feu rouge.
\end{quote}

\begin{quote}
\emph{A reckless driver went through the red light.}
\end{quote}

The suffix \emph{-ard} adds a scornful sense. \emph{Griller un feu} — to run a red light.

\begin{cousinweb}
\emph{chauffard} is built on \emph{chauffeur} \textquotedblleft{}driver\textquotedblright{} with the scornful suffix \emph{-ard}; \emph{chauffer} \textquotedblleft{}to heat\textquotedblright{} is inherited from Latin \emph{calefacere} \textquotedblleft{}to make warm\textquotedblright{} (Lewis \& Short).
\end{cousinweb}

\begin{grammarlens}[title={What you've built}]
One more word for telling, explaining and describing at B1.
\end{grammarlens}

\subsection*{Your turn}
\begin{itemize}
\raggedright
  \item \emph{Say it:} \emph{le chauffard}
  \item \emph{Say it:} \emph{le chauffard}, once more
  \item \emph{Say it:} a sentence of your own with \emph{le chauffard}
  \item \emph{Recall:} say the French for to escape, then the French for the pile-up, then say \emph{le chauffard} again
\end{itemize}

\begin{tcolorbox}[breakable,colback=teal!4,colframe=teal!35!black,title={Before you move on}]
What does le chauffard mean? (\textquotedblleft{}The reckless driver, the road hog\textquotedblright{}.) Say it once more.
\end{tcolorbox}
\section[l'excès de vitesse]{l'excès de vitesse — speeding}

\label{lesson:FR-C297-exces-de-vitesse}



\begin{quote}
Before the new one: say the French for the pile-up, then the French for the reckless driver.
\end{quote}

\subsection*{You'll want to know: l'excès de vitesse}
\textbf{l'excès de vitesse} — \textquotedblleft{}speeding\textquotedblright{}.

\begin{quote}
Il a payé une amende pour excès de vitesse.
\end{quote}

\begin{quote}
\emph{He paid a fine for speeding.}
\end{quote}

\emph{Un excès de vitesse} — a speeding offence.

\begin{cousinweb}
\emph{excès} is a learned borrowing from Latin \emph{excessus} \textquotedblleft{}a going out, a going beyond\textquotedblright{}, from \emph{excēdere} (Lewis \& Short); \emph{vitesse} is built on \emph{vite} \textquotedblleft{}fast\textquotedblright{}, whose origin is not settled.
\end{cousinweb}

\begin{grammarlens}[title={What you've built}]
One more word for telling, explaining and describing at B1.
\end{grammarlens}

\subsection*{Your turn}
\begin{itemize}
\raggedright
  \item \emph{Say it:} \emph{l'excès de vitesse}
  \item \emph{Say it:} \emph{l'excès de vitesse}, once more
  \item \emph{Say it:} a sentence of your own with \emph{l'excès de vitesse}
  \item \emph{Recall:} say the French for the pile-up, then the French for the reckless driver, then say \emph{l'excès de vitesse} again
\end{itemize}

\begin{tcolorbox}[breakable,colback=teal!4,colframe=teal!35!black,title={Before you move on}]
What does l'excès de vitesse mean? (\textquotedblleft{}Speeding\textquotedblright{}.) Say it once more.
\end{tcolorbox}
\section[le constat]{le constat — the accident report (for insurance); a finding}

\label{lesson:FR-C297-constat}



\begin{quote}
Before the new one: say the French for the reckless driver, then the French for speeding.
\end{quote}

\subsection*{You'll want to know: le constat}
\textbf{le constat} — \textquotedblleft{}the accident report (for insurance); a finding\textquotedblright{}.

\begin{quote}
Les deux conducteurs ont rempli un constat.
\end{quote}

\begin{quote}
\emph{The two drivers filled in an accident report.}
\end{quote}

\emph{Le constat amiable} — the joint report that drivers fill in together.

\begin{cousinweb}
\emph{constat} was borrowed from Latin \emph{cōnstat} \textquotedblleft{}it is established, it is agreed\textquotedblright{} (EtymDB), from \emph{cōnstāre} \textquotedblleft{}to stand firm\textquotedblright{}.
\end{cousinweb}

\begin{grammarlens}[title={What you've built}]
One more word for telling, explaining and describing at B1.
\end{grammarlens}

\subsection*{Your turn}
\begin{itemize}
\raggedright
  \item \emph{Say it:} \emph{le constat}
  \item \emph{Say it:} \emph{le constat}, once more
  \item \emph{Say it:} a sentence of your own with \emph{le constat}
  \item \emph{Recall:} say the French for the reckless driver, then the French for speeding, then say \emph{le constat} again
\end{itemize}

\begin{tcolorbox}[breakable,colback=teal!4,colframe=teal!35!black,title={Before you move on}]
What does le constat mean? (\textquotedblleft{}The accident report (for insurance); a finding\textquotedblright{}.) Say it once more.
\end{tcolorbox}
\section[le dérapage]{le dérapage — the skid; a slip-up}

\label{lesson:FR-C297-derapage}



\begin{quote}
Before the new one: say the French for speeding, then the French for the accident report.
\end{quote}

\subsection*{You'll want to know: le dérapage}
\textbf{le dérapage} — \textquotedblleft{}the skid; a slip-up\textquotedblright{}.

\begin{quote}
Sur la route mouillée, le camion a fait un dérapage.
\end{quote}

\begin{quote}
\emph{On the wet road, the lorry went into a skid.}
\end{quote}

Also in politics: \emph{un dérapage verbal} — an unfortunate remark.

\begin{cousinweb}
\emph{dérapage} is built on the verb meaning \textquotedblleft{}to skid\textquotedblright{}, which is usually traced to Provençal and to a Germanic root meaning \textquotedblleft{}to snatch\textquotedblright{}; EtymDB lists no parent.
\end{cousinweb}

\begin{grammarlens}[title={What you've built}]
One more word for telling, explaining and describing at B1.
\end{grammarlens}

\subsection*{Your turn}
\begin{itemize}
\raggedright
  \item \emph{Say it:} \emph{le dérapage}
  \item \emph{Say it:} \emph{le dérapage}, once more
  \item \emph{Say it:} a sentence of your own with \emph{le dérapage}
  \item \emph{Recall:} say the French for speeding, then the French for the accident report, then say \emph{le dérapage} again
\end{itemize}

\begin{tcolorbox}[breakable,colback=teal!4,colframe=teal!35!black,title={Before you move on}]
What does le dérapage mean? (\textquotedblleft{}The skid; a slip-up\textquotedblright{}.) Say it once more.
\end{tcolorbox}
